% Part: second-order-logic
% Chapter: syntax-and-semantics
% Section: expressive-power

\documentclass[../../../include/open-logic-section]{subfiles}

\begin{document}

\olfileid[pt]{sol}{syn}{exp}
\olsection{Poder expressivo}

\begin{explain}
A quantificação sobre variáveis de segunda ordem é responsável por um
enorme aumento no poder expressivo da linguagem em relação ao da lógica de
primeira ordem. A quantificação existencial de segunda ordem nos permite
dizer que existem funções ou relações com certas propriedades. Na lógica de
primeira ordem, a única maneira de fazer isso é especificar um símbolo não
lógico (isto é, !!a{function} ou !!{predicate}) para esse fim. A
quantificação universal de segunda ordem nos permite dizer que todos os
subconjuntos e relações no !!{domain}, ou todas as funções do
!!{domain} para ele, têm uma propriedade. Na lógica de primeira
ordem, só podemos dizer que têm uma propriedade os subconjuntos, relações
ou funções atribuídos a um dos símbolos não lógicos da linguagem. E, ao
dizer que existem subconjuntos, relações ou funções com uma propriedade,
ou que todos eles a têm, também podemos usar quantificação de segunda ordem
para especificar essa propriedade. Isso nos permite definir relações não
definíveis na lógica de primeira ordem e expressar propriedades do domínio
que nela não são expressáveis.
\end{explain}

\begin{defn}
Se $\Struct{M}$ é !!a{structure} para uma linguagem~$\Lang{L}$, uma
relação~$R \subseteq \Domain{M}^2$ é \emph{definível} em~$\Lang{L}$ se
existe alguma !!{formula}~$!A_R(x, y)$ cujas únicas variáveis livres são
$x$ e~$y$, tal que $R(a, b)$ vale (isto é, $\tuple{a,b} \in R$) sse
$\Sat{M}{!A_R(x, y)}[s]$ para $s(x) = a$ e $s(y) = b$.
\end{defn}

\begin{ex}
Na lógica de primeira ordem, podemos definir a relação de identidade
~$\Id{\Domain{M}}$ (isto é, $\Setabs{\tuple{a,a}}{a \in
  \Domain{M}}$) pela fórmula $\eq[x][y]$. Na lógica de segunda ordem,
podemos definir essa relação \emph{sem~$\eq$}. Pois, se $a$ e $b$ são o
mesmo !!{element} de~$\Domain{M}$, então são !!{element}s dos mesmos
subconjuntos de~$\Domain{M}$ (uma vez que conjuntos são determinados por
seus !!{element}s). Reciprocamente, se $a$ e $b$ são diferentes, então não
são !!{element}s dos mesmos subconjuntos: por exemplo, $a \in \{a\}$, mas
$b \notin \{a\}$ se $a \neq b$. Assim, ``serem !!{element}s dos mesmos
subconjuntos de~$\Domain{M}$'' é uma relação que vale para $a$ e $b$ sse
$a = b$. Trata-se de uma relação que pode ser expressa na lógica de segunda
ordem, pois podemos quantificar sobre todos os subconjuntos de~$\Domain{M}$.
Logo, a !!{formula} seguinte define $\Id{\Domain{M}}$:
\[
\lforall[X][(X(x) \liff X(y))]
\]
\end{ex}

\begin{prob}
Mostre que $\lforall[X][(X(x) \lif X(y))]$ (observe:
$\lif$, não~$\liff$!) define $\Id{\Domain{M}}$.
\end{prob}

\begin{ex}
Se $R$ é um !!{predicate} binário, $\Assign{R}{M}$ é uma relação binária
em~$\Domain{M}$. De modo talvez um pouco confuso, usaremos $R$ tanto para o
!!{predicate}~$R$ quanto para a própria relação~$\Assign{R}{M}$. O
\emph{fecho transitivo}~$R^*$ de~$R$ é a relação que vale entre $a$ e $b$
sse, para alguns $c_1$, \dots, $c_k$, valem $R(a,c_1)$, $R(c_1, c_2)$,
\dots, $R(c_k,b)$. Isso inclui o caso $k = 0$, isto é, se $R(a,b)$ vale,
então $R^*(a,b)$ também vale. Isso significa que $R \subseteq R^*$. De
fato, $R^*$ é a menor relação que contém~$R$ e é transitiva. Podemos dizer,
na lógica de segunda ordem, que $X$ é uma relação transitiva que contém~$R$:
\begin{multline*}
  !B_R(X) \ident \lforall[x][\lforall[y][(R(x,y) \lif X(x, y))]] \land {}\\
\lforall[x][\lforall[y][\lforall[z][((X(x,y) \land X(y,z)) \lif X(x,
      z))]]].
\end{multline*}
A primeira conjunção diz que $R \subseteq X$, e a segunda, que $X$ é
transitiva.

Dizer que $X$ é a menor relação desse tipo equivale a dizer que ela própria
está contida em toda relação que contém $R$ e é transitiva. Assim, podemos
definir o fecho transitivo de~$R$ pela !!{formula}
\[
R^*(X) \ident !B_R(X) \land \lforall[Y][(!B_R(Y) \lif
  \lforall[x][\lforall[y][(X(x, y) \lif Y(x,y))]])].
\]
Temos $\Sat{M}{R^*(X)}[s]$ sse $s(X) = R^*$. O fecho transitivo de~$R$ não
pode ser expresso na lógica de primeira ordem.
\end{ex}

\end{document}
